%=====================================================================
\section{Final --- Random reciente}
%=====================================================================

%=====================================================================
\subsection{Ejercicio 1 --- Pulseras BLE del edificio Helíos (AES-CTR, modo degradado)}
%=====================================================================

\begin{consigna}
El edificio Helíos usa pulseras electrónicas para controlar el acceso a las puertas.

Cada pulsera tiene un identificador público $ID$, una clave simétrica $K$ de 128
bits, y un bit \texttt{wearing} que indica si está siendo usada por una persona (1)
o no (0).

Cada 2 segundos, la pulsera transmite vía bluetooth (utilizando BLE broadcast, sin
emparejamiento previo): Pulsera $\to$ aire: $(ID, ctr, C)$, donde $ctr$ es un
contador que empieza en 0 y se incrementa, y $C = \mathrm{AES\text{-}CTR}_k(
\mathit{wearing}\|ctr)$, y $k$ es una clave de 128 bits específica para cada
pulsera.

Los controladores de cada puerta escuchan BLE y reenvían cada mensaje a un
servidor central por un canal considerado seguro: Controlador $\to$ Servidor:
$(door\_id, timestamp, ID, ctr, C)$.

El servidor conoce $K$ para cada $ID$ y descifra de dicho mensaje el valor $C$. Si
de dicho descifrado, $ctr$ coincide con el valor enviado en plano, y el empleado
figura en la base de datos como autorizado para abrir esa puerta, avisa al
controlador que libere la cerradura. Sino, registra el incidente y no permite que
se abra la puerta.

En caso de caída de red (donde los controladores no podrían comunicarse con el
servidor), algunas puertas pasan a modo ``degradado'': si detectan cualquier
mensaje con un ID autorizado recientemente, abren la puerta sin consultar al
servidor.

\begin{enumerate}
  \item[a)] Describir qué verificaciones adicionales podría realizar el servidor
  para detectar usos con los siguientes patrones extraños: intento de abrir dos
  puertas lejanas casi al mismo tiempo, intento de abrir puertas fuera del horario
  normal. (1 pto.)
  \item[b)] Un atacante puede escuchar y transmitir mensajes BLE a su entera
  discreción. Explicar cómo podría modificar o reutilizar mensajes escuchados para
  conseguir abrir una puerta. Proponer una solución. (1 pto.)
  \item[c)] Si un atacante logra forzar el modo degradado, es trivial repetir un
  mensaje reciente y conseguir que se abra la puerta. Proponer alguna mejora que
  mitigue el problema sin introducir nuevos dispositivos o enlaces al sistema.
  (1 pto.)
\end{enumerate}
\end{consigna}

\paragraph{Temas.}
Primer Parte, Clase 2 --- Cifrado (\S Modo CTR: keystream $E_k(\text{nonce}\|i)$,
descifra con la misma función $E$, nunca usa $D_k$, paralelizable, CPA-Secure solo
si no se repite $(k,\text{nonce})$); Primer Parte, Clase 3 --- MACs y Cifrado
Autenticado (\S Maleabilidad de cifrados de flujo/CTR; \S el cifrado por sí solo no
alcanza para integridad; \S MAC: infalsificabilidad, procesa todo el mensaje);
Segunda Parte, Clase 2 --- Autenticación (\S OTP/HOTP: contador sincronizado,
códigos de un solo uso); Segunda Parte, Clase 6 --- Pentesting (áreas candidatas de
hipótesis, aplicable por analogía a los patrones de uso anómalo del inciso a).

\paragraph{Qué necesitás saber.}
\begin{itemize}
  \item \textbf{CTR es un cifrado de flujo, no autenticado.} $c_i = E_k(\text{nonce}
  \|i)\oplus m_i$; el descifrado es $m_i = E_k(\text{nonce}\|i)\oplus c_i$ (nunca se
  usa $D_k$, es la misma función $E$ en ambos sentidos). CTR da \emph{confidencialidad}
  (si no se repite el par clave/nonce), pero \textbf{no} da integridad ni
  autenticación por sí solo.
  \item \textbf{Maleabilidad de cifrados de flujo (Clase 3, Primer Parte).} Si
  $\mathrm{Enc}_k(m)=G(k)\oplus m=c$, entonces $c'=c\oplus x \Rightarrow
  \mathrm{Dec}_k(c')=m\oplus x$: un atacante que conoce (o adivina) un bit del plano
  en una posición conocida puede flippear ese bit del cifrado y, al descifrarse,
  cambia exactamente ese bit del plano, \textbf{sin conocer la clave}.
  \item \textbf{Ningún criptosistema visto es CCA-Secure:} el cifrado está pensado
  para confidencialidad, no para resistir manipulación del texto cifrado; por eso la
  integridad requiere una primitiva aparte, el \textbf{MAC}, que además debe procesar
  la \emph{totalidad} del mensaje.
  \item \textbf{OTP/HOTP (Clase 2, Segunda Parte):} $A$ y $B$ comparten una clave $K$
  y un contador $C$; el valor depende de $(K,C)$, y cada valor sirve una única vez
  (frescura por consumo, sin necesitar conexión).
\end{itemize}

\paragraph{Notas y conexiones.}
El esquema es, otra vez, \emph{confidencialidad sin integridad}: se cifra
$(\mathit{wearing}\|ctr)$ con AES-CTR, pero el mensaje $(ID,ctr,C)$ no lleva ningún
MAC que lo ate. El $ctr$ además viaja \textbf{en claro} junto al cifrado, lo que
facilita comparar y forzar coincidencias. A diferencia del ejercicio análogo de
pulseras BLE resuelto en el final 2025-C2(1) (donde el orden era $CTR\|bit\_wearing$
y no había pregunta sobre patrones de uso anómalo tipo IDS/SIEM, sino sobre estado
degradado con OTP), acá el orden del cifrado es $wearing\|ctr$ (no cambia el
análisis de maleabilidad, solo la posición del bit dentro del bloque cifrado) y el
inciso a) pide explícitamente validaciones de \textbf{correlación entre eventos}
(dos puertas lejanas casi al mismo tiempo, horario anormal), que en el material no
tienen una caja dedicada de IDS/SIEM en las unidades relevadas para este final; se
resuelve igual con el razonamiento de que el servidor puede cruzar información
entre mensajes de distintas puertas y horarios, algo que se desprende directamente
de que el servidor central recibe y procesa \emph{todos} los mensajes de todos los
controladores.

\paragraph{Resolución.}

\textbf{a) Verificaciones adicionales para detectar patrones extraños.}
El servidor central recibe los mensajes de \emph{todas} las puertas (con
$door\_id$ y $timestamp$), por lo que puede cruzar información entre mensajes en
lugar de validar cada uno de forma aislada:
\begin{itemize}
  \item \textbf{Coherencia física (velocidad de desplazamiento):} el servidor
  conoce la ubicación de cada puerta ($door\_id$). Si un mismo $ID$ abre una puerta
  y, en un lapso de tiempo menor al físicamente necesario para recorrer la
  distancia entre ambas, se registra otro intento de apertura de una puerta lejana
  con el mismo $ID$, la solicitud es físicamente imposible y debería rechazarse
  (y generar alerta/incidente), en lugar de autorizar ambas.
  \item \textbf{Ventanas horarias:} asociar a cada $ID$ (o a cada puerta) un rango
  horario habilitado; si el $timestamp$ del mensaje cae fuera de ese rango, denegar
  el acceso aunque el resto de la validación ($ctr$, autorización) sea correcta, y
  registrar el incidente (tal como ya hace el servidor cuando el $ctr$ no coincide
  o el empleado no está autorizado).
  \item En ambos casos, la validación no reemplaza la verificación criptográfica
  del mensaje individual (descifrado + comparación de $ctr$ + autorización): la
  \textbf{complementa}, cruzando el historial de eventos que el servidor ya
  centraliza para detectar patrones anómalos entre mensajes.
\end{itemize}

\textbf{b) Modificación/reutilización de mensajes escuchados.}
Como $C$ se obtiene con AES-CTR (cifrado de flujo), aplica directamente la
maleabilidad de los cifrados de flujo:
\begin{itemize}
  \item \textbf{Replay puro:} el atacante escucha un mensaje válido $(ID, ctr, C)$
  y lo retransmite. El enunciado dice que el servidor compara si el $ctr$
  descifrado coincide con el enviado en plano, pero \textbf{no} dice que lleve un
  registro del último $ctr$ usado por ese $ID$; si no lo lleva, retransmitir el
  mismo mensaje capturado hace que el servidor lo vuelva a aceptar tal cual.
  \item \textbf{Modificación por maleabilidad del bit \texttt{wearing}:} dado que
  $C = E_k(\text{nonce}\|i)\oplus(\mathit{wearing}\|ctr)$, sin conocer $k$ el
  atacante puede flippear el bit del cifrado que corresponde a la posición del bit
  \texttt{wearing} dentro de $C$ (si conoce o adivina esa posición): eso cambia
  exactamente ese bit al descifrar, sin necesitar la clave, igual que el ataque de
  maleabilidad visto en la Clase 3.
  \item \textbf{Solución:} agregar \textbf{integridad} con un MAC sobre el mensaje
  completo $(ID\|ctr\|\mathit{wearing})$, calculado con una clave de MAC (idealmente
  distinta de la de cifrado) y verificado por el servidor antes de aceptar el
  mensaje: cualquier bit flippeado por maleabilidad invalidaría la etiqueta. Además,
  el servidor debe llevar un registro del último $ctr$ aceptado por cada $ID$ y
  solo aceptar $ctr_{\text{nuevo}} > ctr_{\text{último}}$ (monotonicidad/frescura),
  lo que bloquea el replay puro descripto arriba.
\end{itemize}

\textbf{c) Mitigar el replay en modo degradado sin agregar dispositivos ni enlaces.}
En modo degradado la puerta no tiene conexión con el servidor, así que no puede
verificar nada dinámicamente contra un estado remoto; solo puede razonar con la
información que ya escucha localmente (sin sumar hardware ni enlaces nuevos):
\begin{itemize}
  \item Dado que el $ctr$ de cada pulsera es \textbf{estrictamente creciente} (el
  enunciado dice que ``empieza en 0 y se incrementa''), el propio controlador de la
  puerta puede, en modo degradado, \textbf{recordar localmente} (en su memoria) el
  último $ctr$ visto por cada $ID$ que abrió recientemente, sin necesitar consultar
  al servidor ni agregar ningún dispositivo/enlace nuevo.
  \item Con esa información, el controlador degradado deja de aceptar ``cualquier
  mensaje con un ID autorizado recientemente'' y exige, además, que el $ctr$
  recibido sea \textbf{mayor} al último $ctr$ que ya usó ese $ID$ para abrir esa
  puerta: un mensaje reciente repetido (mismo $ctr$) queda rechazado, porque ya fue
  consumido.
  \item Esto es exactamente la misma defensa de monotonicidad de frescura del
  inciso b), aplicada localmente en el controlador (que ya recibe el $ctr$ en
  claro), sin requerir conexión al servidor ni ningún componente adicional.
\end{itemize}

%=====================================================================
\subsection{Ejercicio 2 --- Protocolo de autenticación Cliente--Servidor ($H(K_c\|r)$)}
%=====================================================================

\begin{consigna}
Un sistema de control de acceso utiliza el siguiente protocolo de autenticación
entre un cliente $C$ y un servidor $S$. Cada cliente comparte con el servidor una
clave secreta $K_c$. $H(\cdot)$ es una función de hash libre de colisiones:
\begin{enumerate}
  \item[1.] $C\to S$: $id$
  \item[2.] $S\to C$: $r$ (un número aleatorio de 128 bits)
  \item[3.] $C\to S$: $t = H(K_c\|r)$
  \item[4.] $S$ busca la clave $K_c$ asociada a $id$, recalcula $H(K_c\|r)$ y
  acepta la autenticación si coincide con $t$.
\end{enumerate}

\begin{enumerate}
  \item[a)] Explicar por qué este protocolo es mejor que uno donde el cliente
  enviado solo un hash de la clave secreta. (1 pto.)
  \item[b)] Para evitar la aleatoriedad en el protocolo, un desarrollador decide
  modificar el paso 2 y mantener un contador $ctr$, que transmite en lugar de $r$
  y va incrementando en cada uso. \textquestiondown{}Es esta implementación menos segura que la
  anterior? Justificar. (1 pto.)
  \item[c)] Modificar el protocolo para que exista autenticación mutua, es decir,
  que $C$ también pueda autenticar a $S$ como resultado del intercambio. No
  introducir nuevas claves. (1 pto.)
\end{enumerate}
\end{consigna}

\paragraph{Temas.}
Segunda Parte, Clase 2 --- Autenticación (\S Challenge-Response y EKE: el desafío
$r$ como nonce, ``el nonce evita reutilizar respuestas capturadas''; \S OTP/HOTP:
contador vs.\ aleatoriedad); Primer Parte, Clase 5 --- Protocolos (\S Nonce y
frescura; \S Ataque por reflexión y la necesidad de ``romper la simetría'' de los
mensajes con nonces/identidades distintas por sentido).

\paragraph{Qué necesitás saber.}
\begin{itemize}
  \item \textbf{Challenge-response (Clase 2, Segunda Parte):} el servidor manda un
  challenge único $r$ (nonce); el cliente responde $f(k,r)$ (acá, $H(K_c\|r)$)
  probando que conoce el secreto $K_c$ \textbf{sin enviarlo}. El material remarca
  explícitamente: ``el nonce evita reutilizar respuestas capturadas''.
  \item \textbf{Nonce y frescura (Clase 5, Primer Parte):} un nonce es un número
  aleatorio de un solo uso que aporta frescura a un intercambio y evita ataques de
  \emph{replay}.
  \item \textbf{OTP/HOTP (Clase 2, Segunda Parte):} usar un \textbf{contador}
  sincronizado en vez de un valor aleatorio también da frescura si se exige que sea
  creciente, pero es la base de HOTP, que requiere sincronización de contador entre
  ambas partes (a diferencia de un desafío aleatorio, imprevisible).
  \item \textbf{Ataque por reflexión (Clase 5, Primer Parte):} en un protocolo de
  challenge-response con nonces, un atacante puede abrir una \textbf{sesión
  paralela} y reenviar el mismo desafío recibido para que la víctima ``firme''
  ($f(r)$) y luego reflejar esa respuesta en la sesión original. La defensa que
  remarca el material es \textbf{romper la simetría} de los mensajes (incluir
  identidad/dirección, o usar nonces distintos por sentido).
\end{itemize}

\paragraph{Notas y conexiones.}
Este protocolo es el patrón challenge-response de la Clase 2 (Segunda Parte)
aplicado con un hash en vez de una función $f(k,r)$ genérica. Es prácticamente
idéntico al ejercicio 2 del final 2025-C2(1), con dos diferencias puntuales de
enunciado que no cambian el razonamiento de fondo: acá $r$ se especifica
explícitamente de \textbf{128 bits}, y el inciso a) compara contra un cliente que
enviara ``solo un hash de la clave secreta'' ($H(K_c)$, sin el desafío) en lugar de
preguntar directamente qué pasaría si no estuviese $r$ --- es la misma idea. El
ataque por reflexión de la Clase 5 (Primer Parte) es la advertencia clave para el
inciso c): agregar un segundo desafío simétrico sin distinguir el sentido/rol de
cada uno dejaría abierta la posibilidad de una sesión paralela que refleje la
respuesta.

\paragraph{Resolución.}

\textbf{a) Por qué es mejor que enviar solo $H(K_c)$.}
Si el cliente enviara directamente $H(K_c)$ (sin el desafío $r$), esa respuesta
sería \textbf{siempre el mismo valor fijo} en cada intento de autenticación: un
atacante que la capture una sola vez (escuchando el canal) podría luego
\textbf{reenviarla (replay)} y el servidor la aceptaría como si fuera el cliente
legítimo, sin conocer nunca $K_c$. Con el protocolo dado, en cambio, el servidor
manda en cada autenticación un $r$ \textbf{aleatorio de 128 bits} distinto (un
nonce), por lo que la respuesta $t=H(K_c\|r)$ también es distinta en cada
ejecución: capturar una respuesta pasada no sirve para autenticarse de nuevo, ya
que no hay dos ejecuciones del protocolo con el mismo $r$ (y por lo tanto la misma
$t$) esperado. Es exactamente el rol del nonce en challenge-response: dar frescura
y evitar que se reutilicen respuestas capturadas.

\textbf{b) Contador incremental en vez de $r$ aleatorio: \textquestiondown{}menos seguro?}
Sí, \textbf{es menos seguro} que usar un $r$ aleatorio, aunque un contador
creciente ($ctr$) igualmente aporta algo de frescura si el servidor exige que sea
estrictamente mayor al último usado (evitaría el replay literal del mismo
mensaje). El problema es la \textbf{predictibilidad}: al ser incremental y
conocido (se transmite en claro, igual que $r$), un atacante que observe la
secuencia de valores usados en sucesivas autenticaciones puede \textbf{anticipar
el próximo valor} de $ctr$ antes de que el servidor lo envíe. Esto es análogo a la
distinción entre HOTP (contador, cuya progresión es previsible y requiere
sincronización explícita entre las partes) y un desafío verdaderamente
\textbf{aleatorio} (imprevisible salvo que se conozca la clave). Un valor
aleatorio no permite a un atacante preparar de antemano una estrategia para un
desafío futuro, mientras que un contador sí revela ese futuro con solo observar
el tráfico.

\textbf{c) Autenticación mutua sin introducir nuevas claves.}
Reutilizando la misma $K_c$ compartida, se agrega un desafío en el otro sentido,
rompiendo la simetría (Clase 5, Primer Parte) para evitar un ataque por
reflexión:
\begin{align*}
C \to S &: id,\ r_C\\
S \to C &: r_S,\ H(K_c\|r_C)\\
C \to S &: H(K_c\|r_S)
\end{align*}
El servidor prueba que conoce $K_c$ respondiendo al desafío $r_C$ del cliente (en
el mismo mensaje en que envía su propio desafío $r_S$); el cliente, a su vez,
responde al desafío $r_S$ recibido. Cada parte verifica el hash correspondiente a
\emph{su propio} desafío, y ambos usan la misma $K_c$ y la misma $H$, sin
introducir ninguna clave nueva. Siguiendo la defensa que remarca el material
contra el ataque por reflexión, conviene además que los dos desafíos \textbf{no
sean intercambiables} uno por el otro sin ambigüedad (por ejemplo, atando dentro
del hash la identidad o el rol de quien responde, del estilo
$H(K_c\|r_C\|\text{``servidor''})$ y $H(K_c\|r_S\|\text{``cliente''})$), de modo
que una sesión paralela abierta por un atacante no pueda usarse para que una de
las partes legítimas ``firme'' un desafío que luego se reutilice en la otra
sesión.

%=====================================================================
\subsection{Ejercicio 3 --- CasaIA (IoT hogareño): hipótesis de falla}
%=====================================================================

\begin{consigna}
La empresa CasaIA vende dispositivos inteligentes para hogares (cerraduras
inteligentes, luces, cámaras) que se controlan desde una aplicación móvil y desde
una página web.

Cada usuario crea una cuenta con email y contraseña. La autenticación se hace
contra un API REST centralizado, que devuelve un token de sesión.

Para vincular un dispositivo, el usuario escanea con la app un código QR impreso
en el equipo; este código contiene el identificador del dispositivo y una clave
de emparejamiento.

Los dispositivos se conectan por WiFi a los servidores de CasaIA utilizando un
protocolo propietario sobre TCP, con un canal persistente para recibir comandos
(abrir/cerrar, encender/apagar, etc.).

El sistema permite definir escenarios (por ejemplo, ``modo vacaciones'', ``modo
noche''). Los escenarios se configuran desde la app/web, y el servidor los ejecuta
enviando comandos a los distintos dispositivos según la hora, la geolocalización
del teléfono del usuario, o sensores (movimiento, luz).

El sistema permite otorgar acceso temporal a invitados: el dueño de la casa
ingresa el email de la otra persona, y el sistema le envía un enlace que, al
abrirlo, permite controlar ciertos dispositivos durante una ventana de tiempo
configurable.

En función de la descripción del sistema, establecer 4 hipótesis de falla
plausibles siguiendo la metodología de pentesting. Por cada hipótesis definir una
prueba siguiendo los lineamientos de la metodología. (1 pto c/u.)

\emph{Aclaración:} No utilizar problemas genéricos como hipótesis de falla.
Plantear situaciones y lugares concretos donde podría estar el problema. Por
ejemplo, en lugar de decir ``puede haber un buffer overflow en la lectura del
parámetro x que debe interpretarse como número'', decir ``es probable que haya un
buffer pequeño dado que los números esperados tienen menos de 4 dígitos''.
\end{consigna}

\paragraph{Temas.}
Segunda Parte, Clase 6 --- Pentesting (\S Metodología de hipótesis de falla, Paso 2:
áreas candidatas de hipótesis, comparar con sistemas parecidos; \S Paso 3, la
prueba debe ser repetible y concluyente); Segunda Parte, Clase 3 --- Principios de
Diseño y Vulnerabilidades (\S CSRF; \S XSS, CWE-79); Segunda Parte, Clase 2 ---
Autenticación (\S OTP/HOTP, contador vs.\ aleatorio); glosario (\S OAuth2/tokens,
\S nonce, \S CSRF, \S XSS).

\paragraph{Qué necesitás saber.}
\begin{itemize}
  \item \textbf{Metodología de hipótesis de falla (Clase 6):} el Paso 2 consiste en
  plantear posibles vulnerabilidades examinando implementaciones y mecanismos
  concretos del sistema (\textquestiondown{}está bien implementado el control?, \textquestiondown{}el ambiente
  introduce errores?), y comparando con sistemas parecidos (sistemas similares
  tienen problemas similares). El resultado del Paso 2 es una lista de posibles
  vulnerabilidades asumidas como hipótesis; el Paso 3 (Prueba) exige que la
  verificación de cada hipótesis sea \textbf{repetible y concluyente}.
  \item \textbf{Nonce/token anti-CSRF (Clase 3):} la defensa real contra CSRF es un
  token/nonce único que da frescura y evita replay; el mismo razonamiento aplica a
  cualquier enlace o token de un solo uso (como el enlace de acceso temporal por
  mail de CasaIA).
  \item \textbf{XSS, CWE-79 (Clase 3):} inyectar código que termine ejecutado por
  el browser de la víctima permite robar la sesión, aprovechando la confianza del
  sitio en su sesión/cookie.
  \item \textbf{Tokens de sesión (glosario):} un token de sesión que no tenga
  expiración adecuada, alcance limitado, o que viaje sin protección, puede ser
  robado o reutilizado para suplantar al usuario.
  \item \textbf{OTP/HOTP (Clase 2):} un contador/valor de un solo uso que no se
  invalida correctamente tras su consumo pierde la propiedad de ``one-time'' que
  debería tener.
\end{itemize}

\paragraph{Notas y conexiones.}
La aclaración del enunciado (``no usar problemas genéricos, plantear situaciones y
lugares concretos'') es exactamente el Paso 2 de la metodología de hipótesis de
falla tal como está descripto en el material: no alcanza con decir ``puede haber
una vulnerabilidad de replay'' en abstracto, sino señalar el mecanismo concreto de
CasaIA donde eso es plausible (el QR de vinculación, el canal TCP de comandos, el
enlace de invitado por email, el token de sesión). Cada hipótesis, además, debe
llevar una prueba que sea repetible y concluyente (Paso 3).

\paragraph{Resolución.}
Aplicando la metodología de hipótesis de falla (examinar implementación/mecanismos
concretos + comparar con sistemas parecidos) a cuatro componentes puntuales de
CasaIA (QR de vinculación, canal TCP propietario de comandos, enlace de acceso
temporal por email, token de sesión):

\begin{enumerate}
  \item \textbf{Hipótesis --- el QR de vinculación no vence ni se invalida tras el
  primer uso.} Es probable que el QR impreso en el dispositivo contenga siempre el
  mismo identificador y clave de emparejamiento \emph{estáticos}, sin que el
  servidor los invalide después de la primera vinculación exitosa; si alguien
  fotografía o copia ese QR antes de que el dueño legítimo lo escanee (por ejemplo,
  un QR visible en la góndola de una tienda, o fotografiado por un instalador),
  podría vincular el dispositivo a su propia cuenta antes o después que el dueño
  real. \emph{Prueba (repetible y concluyente):} vincular un dispositivo a una
  cuenta de prueba usando una foto del QR (sin tener el dispositivo físicamente en
  mano) y luego, con el mismo QR, intentar vincularlo desde otra cuenta; si ambas
  vinculaciones son aceptadas (o la segunda no revoca la primera sin aviso al
  dueño), la hipótesis queda confirmada.
  \item \textbf{Hipótesis --- el canal TCP persistente hacia los dispositivos no
  reautentica cada comando individual.} Dado que el protocolo es propietario y el
  canal es persistente (se autentica una vez al abrir la conexión), es probable que
  los comandos posteriores (abrir/cerrar, encender/apagar) que viajan por ese canal
  no lleven ninguna verificación adicional de que provienen de un usuario
  autorizado en ese momento puntual, confiando únicamente en que el socket ya está
  autenticado. \emph{Prueba:} con acceso a la red WiFi del dispositivo (escenario
  de caja gris), capturar el tráfico del canal TCP durante una apertura legítima y
  reenviar (\emph{replay}) ese mismo comando más tarde, o inyectar un comando
  propio en la misma conexión; si el dispositivo lo ejecuta sin re-validar
  autorización, la hipótesis queda confirmada.
  \item \textbf{Hipótesis --- el enlace de acceso temporal por email no tiene
  expiración ni invalidación efectiva.} Es probable que el enlace enviado por
  email para el acceso temporal de invitados sea un token que, una vez generado,
  siga siendo válido más allá de la ventana de tiempo configurada, o que no se
  invalide si el dueño revoca el acceso manualmente antes de tiempo (análogo a un
  token/nonce anti-CSRF sin frescura real); cualquiera que intercepte ese mail
  (reenvío, cuenta de correo comprometida, o el link cacheado en el navegador del
  invitado) podría seguir usándolo después de vencida la ventana. \emph{Prueba:}
  generar un acceso temporal con una ventana corta, dejar que venza (o revocarlo
  manualmente desde la app antes de que venza) y volver a abrir el mismo enlace;
  si el sistema sigue otorgando control sobre los dispositivos, la hipótesis queda
  confirmada.
  \item \textbf{Hipótesis --- el token de sesión queda expuesto a través de un
  campo de texto no sanitizado de la app/web (XSS, CWE-79).} Dado que el sistema
  permite nombrar escenarios (``modo vacaciones'', ``modo noche'') y probablemente
  también nombrar dispositivos o invitados, es probable que alguno de esos campos
  de texto se muestre sin sanitizar en la interfaz web de otro usuario (por
  ejemplo, el dueño ve el nombre que un invitado puso a un escenario compartido);
  si ese campo permite inyectar código, un atacante podría robar el token de
  sesión de quien lo visualiza. \emph{Prueba:} ingresar un script de prueba como
  nombre de un escenario o dispositivo y verificar si se ejecuta al ser
  renderizado en la sesión de otro usuario que visualiza ese mismo escenario/
  dispositivo (por ejemplo, el dueño viendo un escenario creado por un invitado);
  si el script corre y logra leer el token/cookie de sesión, la hipótesis queda
  confirmada.
\end{enumerate}

\begin{ojo}
Como en el análogo del final 2025-C2(1), la metodología de pentesting descripta
en el material de Clase 6 tiene 5 pasos (Recolección $\to$ Hipótesis $\to$ Prueba
$\to$ Generalización $\to$ Eliminación); acá solo se desarrollan explícitamente los
pasos 2 (Hipótesis) y 3 (Prueba) para cada uno de los 4 ítems pedidos, que es lo
que pide puntualmente el enunciado (``4 hipótesis... por cada hipótesis definir
una prueba''). No hay en el material relevado para este final ningún contenido
adicional sobre CasaIA, TCP propietario para IoT, geolocalización de escenarios,
ni protocolos específicos de smart home más allá de los conceptos generales de
CSRF/XSS/tokens/nonce usados arriba: esos detalles del enunciado (protocolo
propietario, escenarios por sensor/geolocalización) no tienen una sección propia
en las clases y se abordan por analogía con los mecanismos generales de
autenticación y control de acceso ya vistos, sin agregar nada externo al material.
\end{ojo}
